\section{Collective Spin-Wave Storage}
\label{sec:collective_spin_wave_storage}

The state stored in the memory is a coherence between the two long-lived states $\ket{g}$ and $\ket{s}$. For a single atom, the local spin-flip operator is defined as~\cite{PsiL-odamIFm_GorshkovEtAl_2007}
\begin{equation}
\hat{\sigma}_{sg}^{(j)}
=
\ket{s_j}\bra{g_j},
\label{eq:single_atom_spin_flip}
\end{equation}
where the superscript $j$ labels the atom within the ensemble. Acting on atom $j$, this operator transfers the population from $\ket{g}$ to $\ket{s}$. In an ensemble memory, the storage interaction does not reveal which atom received the excitation. The corresponding excitation amplitudes must therefore be added coherently over all atoms coupled to the optical mode, producing a delocalized collective excitation~\cite{PsiL-odamIFm_GorshkovEtAl_2007,Lqcwaealo_DuanEtAl_2001}.

For $N$ atoms initially prepared in $\ket{g}$, the general single-excitation spin-wave state may be written as~\cite{PsiL-odamIFm_GorshkovEtAl_2007,Lqcwaealo_DuanEtAl_2001}
\begin{equation}
\ket{1_S}
=
\sum_{j=1}^{N}
c_j
e^{i\Delta\mathbf{k}\cdot\mathbf{r}_j}
\ket{g_1,\ldots,s_j,\ldots,g_N},
\label{eq:single_spin_wave_state}
\end{equation}
where $N$ is the number of participating atoms, $\mathbf{r_j}$ is the position of atom $j$, and $\Delta\mathbf{k}$ is the spin-wave wave vector imposed by the optical fields. The coefficient $c_j$ is the slowly varying excitation amplitude associated with atom $j$ and contains the spatial dependence of the atomic density and light--matter coupling. Since the localized single-excitation states are orthogonal, normalization requires
$\sum_{j=1}^{N}
\left|c_j\right|^2
=
1.$ The absence of which-atom information gives the excitation its collective character, while the position-dependent phase determines how the atomic contributions interfere during retrieval~\cite{PsiL-odamIFm_GorshkovEtAl_2007,Lqcwaealo_DuanEtAl_2001}.

During the write process, the spatial phase difference between the signal and control fields is transferred to the ground-state coherence. For a signal field with wave vector $\mathbf{k_{sig}}^{\mathrm{write}}$ and a control field with wave vector $\mathbf{k_{c}}^{\mathrm{write}}$, the spin-wave wave vector is~\cite{PsiL-odamIFm_GorshkovEtAl_2007,QmfpDp_FleischhauerEtAl_2002}
\begin{equation}
\Delta\mathbf{k}
=
\mathbf{k_{sig}}^{\mathrm{write}}
-
\mathbf{k_{c}}^{\mathrm{write}}.
\label{eq:spin_wave_wavevector}
\end{equation}
The stored excitation therefore forms a spatial grating of ground-state coherence across the ensemble. Atoms at different positions carry different phases, and the preservation of these relative phases determines whether their emitted fields interfere constructively in the selected optical mode during retrieval~\cite{PsiL-odamIFm_GorshkovEtAl_2007,QmfpDp_FleischhauerEtAl_2002}.

The same collective excitation may be represented by the spin-wave creation operator~\cite{PsiL-odamIFm_GorshkovEtAl_2007, Lqcwaealo_DuanEtAl_2001}
\begin{equation}
    \hat{S}_{\Delta\mathbf{k}}^\dagger
    =
    \sum_{j=1}^{N} c_j
    e^{i\Delta\mathbf{k}\cdot\mathbf{r}_j}
    \hat{\sigma}_{sg}^{(j)}.
    \label{eq:spin_wave_operator_discrete}
\end{equation}
Acting on the collective ground state $\ket{G}
    =
    \ket{g_1,\ldots,g_N},
    \label{eq:collective_ground_state}
$ creates the state: 
\begin{equation}
    \ket{1_S}
    =
    \hat{S}_{\Delta\mathbf{k}}^\dagger\ket{G}.
    \label{eq:spin_wave_state_operator}
\end{equation}
For equal coupling to all participating atoms, $c_j=1/\sqrt{N}$, and the operator reduces to the symmetric collective form. In the weak-excitation limit, where the number of stored excitations is much smaller than \(N\), the collective spin-wave mode can be treated approximately as bosonic~\cite{PsiL-odamIFm_GorshkovEtAl_2007}.

For an extended ensemble, the stored excitation can be more naturally described as a continuous spatial coherence. The spin wave contains both an amplitude envelope and a spatial phase grating. Following the standard spin-wave description used in ensemble memories~\cite{PsiL-odamIFm_GorshkovEtAl_2007, QmfpDp_FleischhauerEtAl_2002}, it may be written as
\begin{equation}
    S(\mathbf{r},0)
    =
    f(\mathbf{r})
    e^{i\Delta\mathbf{k}\cdot\mathbf{r}},
    \label{eq:continuous_spin_wave_initial}
\end{equation}
where $S(\mathbf{r},0)$ is the initial ground-state coherence, $f(\mathbf{r})$ is the slowly varying spatial envelope of the stored mode, and $e^{i\Delta\mathbf{k}\cdot\mathbf{r}}$ the phase spatial period written by the signal and control fields. The envelope is set by the optical mode profiles, the atomic density distribution, and the local storage efficiency. During storage, atomic motion can reduce the mode overlap by moving atoms out of the region addressed by the optical fields. It can also reduce the phase-matched collective coherence when atoms move relative to the spin-wave period or when position-dependent energy shifts produce differential phase accumulation~\cite{PsiL-odamIFm_GorshkovEtAl_2007, QmfpDp_FleischhauerEtAl_2002}.
\subsection{Phase-matching condition}
During retrieval, each atom contributes a field amplitude driven by the read control field. The field emitted into a given output direction is the coherent sum of all atomic contributions. Phase-matched contributions add constructively, whereas phase mismatch redirects or suppresses the emission. Retrieval is therefore governed by the spatial phase and envelope of the remaining spin coherence, not only by the amount of microscopic coherence that survives~\cite{PsiL-odamIFm_GorshkovEtAl_2007}.

The phase-matching condition follows from the coherent emission amplitude. During readout, a control field with wave vector \(\mathbf{k}_{\mathrm{c}}^{\mathrm{read}}\) converts the spin wave into an optical field. The emission amplitude into an output mode with wave vector \(\mathbf{k}_{\mathrm{out}}\) contains the factor~\cite{PsiL-odamIFm_GorshkovEtAl_2007}
\begin{equation}
    \mathcal{A}(\mathbf{k}_{\mathrm{out}})
    \propto
    \sum_{j=1}^{N}
    \exp\left[
    i\left(
    \Delta\mathbf{k}
    +
    \mathbf{k}_{\mathrm{c}}^{\mathrm{read}}
    -
    \mathbf{k}_{\mathrm{out}}
    \right)
    \cdot
    \mathbf{r}_j
    \right],
    \label{eq:spin_wave_readout_amplitude}
\end{equation}
where \(\mathcal{A}(\mathbf{k}_{\mathrm{out}})\) is the retrieval amplitude into the output mode. Constructive interference is obtained when the phase is stationary across the ensemble, which gives
$
    \mathbf{k}_{\mathrm{out}}
    =
    \Delta\mathbf{k}
    +
    \mathbf{k}_{\mathrm{c}}^{\mathrm{read}}.
    \label{eq:phase_matching_spin_wave}
$ using Eq.~\eqref{eq:spin_wave_wavevector}, the output wave vector is
\begin{equation}
    \mathbf{k}_{\mathrm{out}}
    =
    \mathbf{k}_{\mathrm{sig}}^{\mathrm{write}}
    -
    \mathbf{k}_{\mathrm{c}}^{\mathrm{write}}
    +
    \mathbf{k}_{\mathrm{c}}^{\mathrm{read}}.
    \label{eq:phase_matching_write_read}
\end{equation}
The retrieved direction is therefore fixed by the spin-wave wave vector written during storage and by the geometry of the read control field~\cite{PsiL-odamIFm_GorshkovEtAl_2007}.

The spin wave is sensitive to atomic motion because its phase is evaluated at the atomic positions. If atom \(j\) moves from \(\mathbf{r}_j(0)\) to \(\mathbf{r}_j(t)\), the factor \(e^{i\Delta\mathbf{k}\cdot\mathbf{r}_j}\) changes. A common displacement of all atoms gives only a global phase, while random thermal motion changes the relative phases across the ensemble. The time-dependent spin-wave state may be written schematically as~\cite{PsiL-odamIFm_GorshkovEtAl_2007}
\begin{equation}
    \ket{1_S(t)}
    =
    \frac{1}{\sqrt{N}}
    \sum_{j=1}^{N}
    e^{i\Delta\mathbf{k}\cdot\mathbf{r}_j(t)}
    e^{-i\phi_j(t)}
    \ket{g_1,\ldots,s_j,\ldots,g_N},
    \label{eq:time_dependent_spin_wave}
\end{equation}
where \(\phi_j(t)\) denotes any additional phase accumulated by atom \(j\) during storage. Randomization of the relative phases reduces the collective mode that emits into the desired optical field, even if local ground-state coherence remains~\cite{PsiL-odamIFm_GorshkovEtAl_2007}.


\subsection{Maxwell--Bloch Description of Light Storage and Retrieval}

The collective spin-wave description introduced above specifies the spatial form of the stored excitation but does not describe its conversion to and from the optical field. To represent this dynamics, the one-dimensional Maxwell--Bloch formulation developed by Gorshkov \emph{et al.} is adopted in this subsection~\cite{PsiL-odamIFm_GorshkovEtAl_2007}. In particular, the definitions, normalization conventions, and coupled equations presented below follow their treatment of storage and retrieval in an optically dense $\Lambda$-type medium.

Within this formulation, the system is described by the slowly varying signal-field amplitude $\mathcal{E}(z,t)$, the optical polarization $P(z,t)$, and the spin wave $S(z,t)$. Following the definitions of Gorshkov \emph{et al.}~\cite{PsiL-odamIFm_GorshkovEtAl_2007}, the collective atomic variables are written as
\begin{equation}
P(z,t)
=
\sqrt{N},\sigma_{ge}(z,t),
\qquad
S(z,t)
=
\sqrt{N},\sigma_{gs}(z,t),
\label{eq}
\end{equation}
where $\sigma_{ge}(z,t)$ is the slowly varying optical coherence on the $\ket{g}\leftrightarrow\ket{e}$ transition and $\sigma_{gs}(z,t)$ is the ground-state coherence on the $\ket{g}\leftrightarrow\ket{s}$ transition. The polarization $P$ couples directly to the signal field and decays through the excited state, whereas $S$ describes the collective ground-state excitation.

In the one-dimensional treatment, the signal is assumed to remain fixed to a transverse mode throughout the ensemble, such that diffraction and coupling to other transverse modes can be neglected over the interaction length, as assumed by Gorshkov \emph{et al.}~\cite{PsiL-odamIFm_GorshkovEtAl_2007}. The control field is taken to be approximately uniform across this mode, ensuring that the light--matter coupling depends only on the longitudinal position. The transverse overlap of the fields with the atomic ensemble is then incorporated into an effective coupling constant, and the dynamical evolution is described solely along the propagation coordinate $z$.
In a frame co-moving with the signal field and using the dimensionles svariables, the Maxwell--Bloch equations take the form ~\cite{PsiL-odamIFm_GorshkovEtAl_2007}:
\begin{equation}
    \partial_{\tilde{z}}\mathcal{E}
    =
    i\sqrt{d}\,P,
    \label{eq:gorshkov_field}
\end{equation}
\begin{equation}
    \partial_{\tilde{t}}P
    =
    -\left(1+i\tilde{\Delta}\right)P
    +
    i\sqrt{d}\,\mathcal{E}
    +
    i\tilde{\Omega}(\tilde{t})S,
    \label{eq:gorshkov_polarization}
\end{equation}
and
\begin{equation}
    \partial_{\tilde{t}}S
    =
    i\tilde{\Omega}^{*}(\tilde{t})P.
    \label{eq:gorshkov_spin_wave}
\end{equation}

Here, $\tilde{z}$ denotes the dimensionless longitudinal coordinate and $\tilde{t}$ the dimensionless retarded time. The remaining dimensionless parameters are
$
\tilde{\Delta}=\frac{\Delta}{\gamma},
\qquad
\tilde{\Omega}=\frac{\Omega}{\gamma},
$
where $\gamma$ is the optical-polarization decay rate, $\Delta$ is the one-photon detuning, and $\Omega$ is the control-field Rabi frequency~\cite{PsiL-odamIFm_GorshkovEtAl_2007}.

The coupling between the ensemble and the selected optical mode is governed by the resonant optical depth $d$. For a fixed atomic transition, $d$ is determined by the atomic density and the longitudinal extent of the medium. An increase in $d$ strengthens the collectively enhanced coupling to the phase-matched optical mode relative to the decay of the optical polarization into other electromagnetic modes~\cite{PsiL-odamIFm_GorshkovEtAl_2007}.

A central result derived by Gorshkov \emph{et al.} is that the retrieval efficiency from a prescribed spin wave is determined by $d$ and by the spatial profile of that spin wave~\cite{PsiL-odamIFm_GorshkovEtAl_2007}. Provided that retrieval is complete and no excitation remains in $P$ or $S$, the integrated retrieval efficiency is independent of the control-field envelope and the one-photon detuning. These parameters determine the temporal profile of the retrieved optical field, but not the total retrieval efficiency associated with the initial spin-wave mode.

%\subsection{B field}
%Magnetic-field gradients provide a direct source of such dephasing. A spatially varying field changes the local ground-state splitting according to
%\begin{equation}
%    \omega_{sg}(\mathbf{r})
%    =
%    \omega_{sg}^{(0)}
%    +
%    \delta\omega_{sg}(\mathbf{r}),
%    \label{eq:position_dependent_ground_state_splitting}
%\end{equation}
%where \(\omega_{sg}^{(0)}\) is the spatially uniform splitting and \(\delta\omega_{sg}(\mathbf{r})\) is the position-dependent Zeeman shift. The relative phase accumulated by atom \(j\) is then
%\begin{equation}
%    \phi_j(t)
%    =
%    \int_{0}^{t}
%    \delta\omega_{sg}\!\left(\mathbf{r}_j(t')\right)
%    \,dt',
%    \label{eq:magnetic_phase_accumulation}
%\end{equation}
%where \(t'\) is the integration variable. Spatial variation of \(\phi_j(t)\) changes the spin-wave phase pattern and reduces retrieval into the phase-matched optical mode~\cite{PsiL-odamIFm_GorshkovEtAl_2007,Apgteitiav_FinkelsteinEtAl_2023}.
%
%The stored mode may therefore be represented as
%\begin{equation}
%    S(\mathbf{r},t)
%    =
%    f(\mathbf{r},t)
%    e^{i\Delta\mathbf{k}\cdot\mathbf{r}}
%    e^{-i\phi(\mathbf{r},t)},
%    \label{eq:general_spin_wave_mode}
%\end{equation}
%where \(f(\mathbf{r},t)\) is the time-dependent amplitude envelope and \(\phi(\mathbf{r},t)\) is the accumulated position-dependent phase~\cite{PsiL-odamIFm_GorshkovEtAl_2007}. Optical readout is controlled by the overlap between this evolved spin wave and the collective mode satisfying the retrieval phase-matching condition.
